\documentclass[UTF8,zihao=-4,fontset=windows]{ctexart}
\usepackage{amsmath,amssymb}
\usepackage[a4paper,margin=2.15cm]{geometry}
\usepackage{enumitem}
\setlist[enumerate]{itemsep=0.45em,topsep=0.25em,leftmargin=2.2em}
\setlength{\parindent}{0pt}
\setlength{\parskip}{0.34em}

\newcommand{\Prob}{\mathbb P}
\newcommand{\EX}{\mathbb E}
\newcommand{\DD}{\operatorname D}
\newcommand{\iid}{\overset{\mathrm{i.i.d.}}{\sim}}
\newcommand{\Var}{\operatorname{Var}}
\newcommand{\Cov}{\operatorname{Cov}}

\title{\vspace{-1.2cm}随机数学与统计期末考试\\{\large 2020 年秋季回忆版复原}}
\date{}

\begin{document}
\maketitle
\vspace{-1.6em}

说明：本卷根据 2020 年 12 月 30 日期末考试回忆稿复原。个别回忆稿中缺失的小问已按相同题型作最小补全。

\section*{一、}
设 $X,Y,Z \iid U(0,1)$。
\begin{enumerate}
  \item 求 $X+Y+Z$ 的分布函数。
  \item 求
  \[
    \EX\left(X\mid X+Y=\frac12\right),\qquad
    \EX\left(X\mid X+Y\le \frac12\right).
  \]
  \item 求 $-\ln X$ 的矩母函数，并证明
  \[
    U=(XY)^Z\sim U(0,1).
  \]
\end{enumerate}

\section*{二、}
设随机变量 $X$ 的密度函数为
\[
  f_X(x)=
  \begin{cases}
    3x^2, & 0\le x\le 1,\\
    0, & \text{其他}.
  \end{cases}
\]
给定 $X=x$ 时，$Y$ 服从参数为 $x$ 的几何分布，即
\[
  \Prob\{Y=k\mid X=x\}=x(1-x)^{k-1},\qquad k=1,2,\ldots.
\]
\begin{enumerate}
  \item 求 $\EX Y$ 和 $\DD Y$。
  \item 求
  \[
    \Prob\left\{X\le \frac12 \EX(XY)\right\}.
  \]
  \item 求 $Y$ 的分布。
\end{enumerate}

\section*{三、}
设 $(X,Y)$ 在区域
\[
  D=\{(x,y):x>0,\ y>0,\ x+y\le 1\}
\]
上服从均匀分布。
\begin{enumerate}
  \item 求 $\EX\left(Y\mid X=\frac12\right)$。
  \item 求
  \[
    \EX(Y^2\mid X),\qquad \EX(Y^2\mid X<Y).
  \]
  \item 求 $X+Y$ 的密度函数。
  \item 判断 $X$ 与 $Y-\EX(Y\mid X)$ 是否相关、是否独立，并说明理由。
\end{enumerate}

\section*{四、}
设 $\{B_t:t\ge 0\}$ 为标准布朗运动，$B_0=0$。
\begin{enumerate}
  \item 判断 $\{B_{2t}-B_t:t\ge 0\}$ 是否为标准布朗运动，并说明理由。
  \item 判断随机向量
  \[
    \left(B_2-\frac12B_1,\ B_1,\ B_1+aB_2\right)^{\mathsf T}
  \]
  是否服从三维正态分布；若服从，求其均值向量和协方差矩阵。再求常数 $a$，使 $B_1$ 与 $B_1+aB_2$ 独立。
  \item 求 $\Cov(B_2-B_1,\ B_3-B_2)$，并判断二者是否独立。
  \item 求
  \[
    \Prob\{B_1+B_2<3\mid B_3=4\}.
  \]
\end{enumerate}

\section*{五、}
设随机变量 $X$ 的密度函数为
\[
  f_X(x;\theta)=
  \begin{cases}
    e^{-(x-\theta)}, & x>\theta,\quad \theta>0,\\
    0, & \text{其他},
  \end{cases}
\]
$X_1,\ldots,X_n$ 为来自该总体的样本。
\begin{enumerate}
  \item 求 $\theta$ 的矩估计 $\hat\theta_M$ 和最大似然估计 $\hat\theta_{\mathrm{MLE}}$。
  \item 判断 $\hat\theta_{\mathrm{MLE}}$ 是否无偏；若不是，求其偏差。
  \item 若先验分布为 $\theta\sim \operatorname{Exp}(1)$，求 $\theta$ 的后验分布函数。
  \item 设 $N$ 为与样本独立的正整数值随机变量，且 $\EX(1/N)<\infty$。记
  \[
    X_{(1)}=\min\{X_1,\ldots,X_N\}.
  \]
  证明
  \[
    \EX X_{(1)}=\theta+\EX\left(\frac1N\right).
  \]
\end{enumerate}

\section*{六、}
设随机变量 $X$ 的密度函数为
\[
  f_X(x;\mu,\sigma)=
  \begin{cases}
    \dfrac{1}{\sigma}\sqrt{\dfrac2\pi}\,
    \exp\left\{-\dfrac{(x-\mu)^2}{2\sigma^2}\right\},
    & x>\mu,\\
    0, & \text{其他},
  \end{cases}
\]
其中 $\mu\in\mathbb R,\ \sigma>0$，$X_1,\ldots,X_n$ 为来自该总体的样本。
\begin{enumerate}
  \item 若 $Z\sim N(0,1)$，证明 $X$ 与 $\mu+\sigma |Z|$ 同分布。
  \item 求 $(\mu,\sigma)$ 的充分统计量。
  \item 求 $\hat\mu_{\mathrm{MLE}}$ 和 $\hat\sigma^2_{\mathrm{MLE}}$。
  \item 当 $\mu$ 已知时，对
  \[
    H_0:\sigma=1\quad \text{vs.}\quad H_1:\sigma=2
  \]
  作似然比检验，给出水平为 $\alpha$ 的拒绝域。
\end{enumerate}

\section*{附加题}
设 $X_1,\ldots,X_n \iid U(0,1)$，令
\[
  S_n=\left(\prod_{i=1}^n X_i\right)^{-1/n}.
\]
证明
\[
  \sqrt n\,(S_n-e)\xrightarrow{D}N(0,e^2).
\]

\vspace{1em}
{\small 注：原回忆稿第 1 题第 3 小问指数、第 3 题第 3 小问和第 4 题第 3 小问不完整。本复原卷按相邻题意作了最小补全；第 1 题第 3 小问按可成立形式复原为 $U=(XY)^Z$；第 6 题第 4 小问为使检验问题明确，补充“$\mu$ 已知”。}

\end{document}
